The absolute rule used in the main proof requires a fresh-coordinate
bound independent of the dimension of the box. This uniformity
requires more than positive counting. Rational and symbol
calculus in Appendices~\ref{prelim:section}
and~\ref{prelim:step-drop-section} supports the shift comparison of
Appendix~\ref{shift:section}; positive counting in
Appendix~\ref{counting:section} then gives an absolute increment in
fixed dimension. Its extension to arbitrary dimension is the final
relative-lifting argument.

Sampling supplies the connection. After scalar transfer on ordinary
affine samples in Appendix~\ref{sampling:section},
Appendix~\ref{sampling:constrained-section} constructs one ensemble
carrying the required polynomial lifts. It has two distinct uses.
Cube comparison in Appendix~\ref{cube:section} transfers detection
and yields finite linear models. Those models need not have positive
coefficients, so they alone do not return a positive score. The
nonnegative comparison weights of Appendix~\ref{cube:transfer-section}
provide that return and recover the required integer slots.

Appendix~\ref{lifting:section} uses detection and approximation to
replace local orbits by an ambient orbit with an exact marked
projection. It combines this globalization with positive score
transfer to prove relative lifting with rank at most $d+d_0$.
Applying it to the
constant old patch, of rank zero, gives the dimension-independent
fresh rank in Lemma~\ref{ap:absolute}. This argument permits the
stated late costs in the returned patch. The main triangular return
has the separate task of retaining the old determining equations
with width losses that can be iterated. Table~\ref{ap:interface-table}
records the precise analytic inputs and their proofs.

Throughout the appendices, degree parameters are fixed, while dimensions
are included in the stated logarithmic budgets. The polynomial notation
of Section~\ref{ap:setup} is repeated where needed. The refined
width estimates required for the triangular increment are proved in
the main text; they are not hypotheses of these analytic results.

\begin{table}[htbp]
\centering\small
\begin{tabular}{@{}p{.23\linewidth}p{.42\linewidth}p{.27\linewidth}@{}}
\toprule
Input & Result used & Full statement and proof \\
\midrule
A bounded function with large cube norm
& Correlation with a controlled niltest; interval input and local
box/product-cyclic reductions
& Theorem~\ref{prelim:inverse} \\[4pt]
A biased top-degree niltest
& Slow--structured--rational symbol decomposition, with all parameter
dimensions included in the budget
& Lemma~\ref{prelim:step-drop} \\[4pt]
Positive comparison against degree-$d$ tests
& Uniform comparison against shifted degree-$(d-1)$ tests outside a
quantitatively small exceptional set
& Theorem~\ref{shift:required} \\[4pt]
A progression-free input on fixed-dimensional prime boxes
& An absolute increment and conversion between niltests and patches
& Proposition~\ref{counting:absolute};
Lemma~\ref{counting:conversion} \\[4pt]
High formal rank, with structural choices fixed before precision
& One ensemble satisfying geometry, detection and transfer for the
prescribed inputs and accuracies
& Proposition~\ref{sampling:constrained}; proofs in
Appendices~\ref{cube:section} and~\ref{cube:transfer-section} \\[4pt]
A fixed strong projection and exact marked polynomial
& Globalization preserving the marked identity with early output bounds
& Proposition~\ref{lifting:candidate-globalization} \\[4pt]
An absolute rule and an old patch of rank $d$
& A relative increment with rank at most $d+d_0$, including its
separate impossibility clause
& Proposition~\ref{counting:relative-target}; proof in
Theorem~\ref{lifting:relative-theorem} \\
\bottomrule
\end{tabular}
\caption{Analytic inputs to the main proof. Each reference specifies the
complete hypotheses and quantitative bounds; the table records the
object passed between the statements.}
\label{ap:interface-table}
\end{table}
